\documentclass[11pt,a4paper]{article}
\usepackage[margin=27mm]{geometry}
\usepackage{fontspec,unicode-math}
\setmainfont{STIXGeneral}
\setmathfont{STIX Two Math}
\usepackage{amsmath,amsthm,booktabs,hyperref}
\newtheorem{proposition}{Proposition}
\newtheorem{lemma}[proposition]{Lemma}
\newtheorem{corollary}[proposition]{Corollary}
\newcommand{\Aut}{\operatorname{Aut}}
\newcommand{\Inn}{\operatorname{Inn}}
\newcommand{\Out}{\operatorname{Out}}
\newcommand{\F}{\mathbb F}
\title{Borel subgroups and duality in coset twisting}
\author{Independent correction note for the N006 source contract}
\date{11 October 2026}
\begin{document}
\maketitle
\begin{abstract}
We check several literal statements in the 28 May 2026 arXiv revision of
Piedade--Tranchida, \emph{The Geometry of Wreath and Semi-Direct Products}.
A thin rank-one wreath example gives the incidence graph $K_{2,8}$.
Self-dual polytopes of type $\{2,2\}$ and $\{2,2,2\}$ disprove, respectively,
the unqualified automorphism embedding and the universal nonlinear-diagram
conclusion. We also derive the actual wreath parabolics, separate two missing
hypotheses in the preliminary parabolic lemma, and exhibit an outer
automorphism quotient with no homomorphic section. These are bounded
corrections, not a review or recovery of the whole paper.
\end{abstract}

\section{The conventions that matter}
All page references refer to the fixed 22-page source~\cite{source}.
Physical and printed page numbers agree. Its Section 2.1
allows a group action on a geometry through its image in the automorphism
group (p.~3). Thus a flag-transitive coset model need not be faithful. Write
\[
 A_J=\bigcap_{j\in J}A_j,\qquad A_\varnothing=A.
\]
In particular, $C=A_{I_\alpha}$ is the Borel subgroup. This intersection
notation is different from the subgroup generated by a specified subset of
distinguished generators. Definition 3.10 uses the intersection notation.
The string $C$-group definition on p.~5 imposes commutation on
nonconsecutive generators but does not require consecutive products to have
order greater than two. We shall prove directly that the examples using
such products are admissible abstract regular polytopes.

\section{Wreath parabolics and thinness}
Let $B$ act on a finite set $\Omega$. The product geometry has group
$\bar A=\prod_{\omega\in\Omega}{}^\omega A$ and maximal parabolic at
$(j,\omega)$ obtained by replacing the $\omega$-factor by $A_j$.
An orbit of types is $O\times\{j\}$, where $O$ is a $B$-orbit on $\Omega$.
Choose a representative $\lambda_{O,j}\in O$ for each such orbit.
For $i\in I_\beta$ put
\[
 J_i(\omega)=\{j\in I_\alpha:\ 
     \omega\notin B_i\lambda_{O(\omega),j}\}.
\]
\begin{proposition}\label{p:factor}
The parabolic prescribed by Definition 3.10 for a type $i\in I_\beta$ is
\begin{equation}\label{eq:general}
 G_i=\left(\prod_{\omega\in\Omega}{}^\omega A_{J_i(\omega)}\right)\rtimes B_i.
\end{equation}
If the representatives $\lambda_{O,j}=\lambda_O$ are independent of $j$, set
$U_i=\bigcup_O B_i\lambda_O$. Then this becomes
\begin{equation}\label{eq:coherent}
 G_i=\left(\prod_{\omega\in U_i}{}^\omega A\ \times
               \prod_{\omega\notin U_i}{}^\omega C\right)\rtimes B_i.
\end{equation}
\end{proposition}
\begin{proof}
Definition 3.10 intersects the product parabolics indexed by the complement
of each $B_i$-orbit of the chosen representative. At coordinate $\omega$,
these are exactly the restrictions $A_j$ with $j\in J_i(\omega)$.
Intersections in a direct product are coordinatewise, proving
\eqref{eq:general}. With coherent representatives, $J_i(\omega)$ is empty
inside $U_i$ and is all of $I_\alpha$ outside it. This gives
\eqref{eq:coherent}. Each indicated product is $B_i$-invariant, as required
for the displayed semidirect product.
\end{proof}
Proposition 4.3 on p.~15 instead prints full $A$-factors on the complement
of $U_i$, and does not retain the Borel factors. Even when $C=1$, the
support of the full factors should be $U_i$, not its complement. When
representatives depend on $j$, formula \eqref{eq:general} is needed.
The worked symmetric-group example on p.~17 does use the support $U_i$,
consistent with Definition 3.10 and \eqref{eq:coherent}.

\begin{proposition}\label{p:thin}
The thinness conclusion of Theorem 4.4(d), under its literal stated
hypotheses, fails for rank-one thin factors. There is a twisting with
group order $48$ and incidence graph $K_{2,8}$.
\end{proposition}
\begin{proof}
Take $\alpha=(C_2,(1))$ and $\beta=(S_3,(C_3))$, with their single types.
Each geometry consists of two elements, so it is thin and flag-transitive;
rank-one residual connectedness is vacuous under the source definition.
Let $S_3$ act naturally on $\Omega=\{1,2,3\}$. Both $S_3$ and its Borel
$C_3$ are transitive on $\Omega$. In the rank-one orbit-intersection
condition, the two subsets of $I_\beta$ therefore both yield $\Omega$.
All four required intersection equalities hold.

Now $\bar A=C_2^3$ and $G=C_2^3\rtimes S_3$. There is one type orbit in
the product. Definition 3.10 gives the two maximal parabolics
\[
 P=C_2^3\rtimes C_3,\qquad Q=S_3,
 \qquad P\cap Q=C_3.
\]
Indeed the $B_i$-orbit is all three coordinates, so its complement imposes
no restrictions on $\bar A$; the parabolic for the type orbit imposes all
three restrictions and leaves only $Q$. Thus $|P|=24$, $|Q|=6$ and
$|P\cap Q|=3$. The product $PQ$ has $48$ elements and equals $G$.
Every left coset of $P$ consequently meets every left coset of $Q$.
There are two cosets of $P$ and eight cosets of $Q$, proving the asserted
incidence graph. The two panel sizes are $[P:P\cap Q]=8$ and
$[Q:P\cap Q]=2$, so the geometry is not thin.
\end{proof}
Theorem 3.16(e) on p.~12 explicitly assumes trivial Borel subgroups for
thinness, whereas Theorem 4.4(d) omits that assumption. Replacing $\beta$
by its faithful two-point model does not preserve this example's action on
$\Omega$: the normal subgroup $C_3$ acts nontrivially there, so the action
does not factor through $S_3/C_3$. The paragraph following Theorem 3.16
therefore does not repair this particular application. This does not
disprove Corollary 3.17 when its geometries are modeled by their faithful
automorphism groups; the interpretation of the acting models must be stated.

\section{Self-dual polytopes with all adjacent labels two}
\begin{lemma}\label{l:digon}
For every $n\geq2$ there is a self-dual abstract regular $n$-polytope with
automorphism group $V=\F_2^n$, distinguished generators the coordinate
vectors $e_0,\ldots,e_{n-1}$, and type $\{2,\ldots,2\}$.
Its associated coset geometry has two objects of each type, with every
two objects of different types incident.
\end{lemma}
\begin{proof}
Adjoin a least face of rank $-1$ and a greatest face of rank $n$ to two
faces $(i,0),(i,1)$ of each rank $0\leq i<n$. Declare every face of lower
rank smaller than every face of higher rank; the two faces at one rank
are incomparable. Every maximal chain has one face of each rank. An
interval whose endpoints differ by two ranks has exactly the two faces
of the intermediate rank, so the diamond condition holds. A flag is a
choice at some of the ranks. Its completing chambers are all independent
choices at the remaining ranks, and any two are joined by changing one
choice at a time. This also proves strong flag connectivity while keeping
any common faces fixed, for the polytope and each of its sections.

A rank-preserving automorphism can independently swap the two faces of
each proper rank, and has no other freedom. The full group is therefore
$V$, acting regularly on the $2^n$ chambers. The map
$(i,\varepsilon)\mapsto(n-1-i,\varepsilon)$, interchanging least and
greatest faces, is an order-reversing duality fixing the chamber of zero
choices. It induces the involutory automorphism
$\tau(e_i)=e_{n-1-i}$ of $V$.

For completeness, the subgroup generated by $\{e_i:i\in S\}$ consists
exactly of vectors supported on $S$. Two such subgroups intersect in the
one supported on the intersection of their supports. This proves the full
$C$-group intersection property; all distinct generator products have
order two, so the string relations hold. The cosets of
$V_i=\langle e_j:j\ne i\rangle$ are specified by the $i$th coordinate.
Different coordinate conditions are compatible, so the coset geometry is
exactly the incidence geometry just described. It is thin, flag-transitive
and residually connected: every residue of rank at least two is a complete
multipartite graph with two vertices in each remaining part.
\end{proof}
This explicit construction verifies admissibility without assuming that
the source uses a restrictive convention excluding digons. Its definitions
on pp.~3 and 5--6 contain no such restriction.

\begin{proposition}\label{p:embedding}
The conclusion of Theorem 5.2(d) fails for the rank-three polytope of
Lemma~\ref{l:digon}. More generally, if $\tau\in\Aut(A)$ is an outer
involution, the canonical conjugation map has kernel $Z(A)$ and gives
\[
 (A\rtimes\langle\tau\rangle)/Z(A)
       \ \cong\ \langle\Inn(A),\tau\rangle\ \leq\ \Aut(A).
\]
Consequently $Z(A)=1$ is a sufficient extra hypothesis for the claimed
canonical embedding.
\end{proposition}
\begin{proof}
For $A=\F_2^3$, reversal swaps $e_0,e_2$ and fixes $e_1$. It is nontrivial,
and hence outer because $A$ is abelian. The semidirect product has order
$16$, whereas
\[
 |\Aut(A)|=|\operatorname{GL}_3(2)|=(8-1)(8-2)(8-4)=168.
\]
Since $16$ does not divide $168$, no embedding exists, even an unspecified
one. In general the map $(a,\tau^\varepsilon)\mapsto
\operatorname{inn}(a)\tau^\varepsilon$ is a homomorphism onto the displayed
subgroup. An element of its kernel with $\varepsilon=1$ would make $\tau$
inner, which is excluded. With $\varepsilon=0$, the kernel condition is
exactly $a\in Z(A)$. The center is $\tau$-invariant, so the quotient is
well-defined and the assertion follows.
\end{proof}

\begin{proposition}\label{p:nonlinear}
The nonlinear-diagram assertion in Theorem 5.2(b) fails for the rank-four
polytope of Lemma~\ref{l:digon}. Its group $H=\F_2^4\rtimes\langle\tau\rangle$
has no rank-three regular hypertope with a nonlinear diagram, irrespective
of the chosen distinguished generators.
\end{proposition}
\begin{proof}
Write elements as $(v,\varepsilon)$, with multiplication
\[
 (v,\varepsilon)(w,\delta)
       =(v+\tau^\varepsilon w,\varepsilon+\delta).
\]
All involutions in $V\times\{0\}$ commute. An element $(v,1)$ is an
involution exactly when $v+\tau v=0$, that is, $v$ is fixed by $\tau$.
Two such involutions commute, since their products in either order equal
$(v+w,0)$. Thus the involutions split into two sets, each with pairwise
commuting elements. Every triple of involutions contains a commuting pair.

We recall why this is a geometric obstruction. In a thin residually
connected flag-transitive geometry, the chamber graph is connected.
This follows by induction on rank: chambers through a fixed object are
connected by induction in its residue, and a path in the incidence graph
joins these connected sets; rank one starts with its two chambers.
A type-preserving automorphism fixing one chamber fixes each unique
adjacent chamber by thinness, hence all chambers by connectedness, and
therefore every object. The group consequently acts regularly on
chambers. The unique automorphism sending a base chamber to its
$i$-adjacent chamber is an involution $\rho_i$. These involutions generate
the group because the chamber graph is connected. For two types $i,j$,
the rank-two residue has chamber graph the alternating cycle generated by
$\rho_i,\rho_j$; its gonality is the order of $\rho_i\rho_j$. In particular,
commuting distinct generators give gonality two and no diagram edge.

In rank three, every diagram with at least one missing edge is linear
after ordering its three vertices along a path. A nonlinear diagram
therefore needs three pairwise noncommuting distinguished involutions,
which $H$ does not have. The rank prescribed by Theorem 5.2(b) for $n=4$
is three, proving the contradiction. The group has order $32$ and exponent
four, but neither a small-group classification nor a separate order bound
is used in this argument.
\end{proof}
The existence portions of Theorem 5.2(b),(c) and the reversal in (a) are
not refuted by these examples. Proposition~\ref{p:nonlinear} addresses
the additional universal nonlinear choice, and Proposition~\ref{p:embedding}
addresses the assertion of an embedding for every outer reversal.

\section{Two distinct problems with the preliminary parabolic lemma}
The source defines $G^i=G_{I\setminus\{i\}}$ and
$G^J=\langle G^i:i\in J\rangle$, so $G^\varnothing=1$. Lemma 2.4(c),(d)
on pp.~4--5 assert the complementary equalities between $G_J$ and
$G^{I\setminus J}$ for every subset, including the extremes. At $J=I$
in (c), or $J=\varnothing$ in (d), this forces $G_I=1$, which residual
connectedness does not force for an unfaithful coset model.

\begin{proposition}\label{p:lemma}
Under the literal conventions of Lemma 2.4, residual connectedness does
not imply (c)--(e). Conversely, (e) alone need not imply residual
connectedness, even for a thin flag-transitive coset geometry with trivial
Borel subgroup.
\end{proposition}
\begin{proof}
For the first assertion take $G=\langle x,y,z\rangle\cong C_2^3$ and
$G_0=\langle y,z\rangle$, $G_1=\langle x,z\rangle$. The coset geometry is
$K_{2,2}$, so it is thin, flag-transitive and residually connected. Its
Borel is $\langle z\rangle\ne1$, disproving (c),(d) at the extremes.
Moreover
\[
 G^{\{0\}}\cap G^{\{1\}}=G_1\cap G_0=\langle z\rangle
       \ne G^\varnothing,
\]
so (e) also fails.

For the converse use the same $G$, but $G_0=\langle x\rangle$ and
$G_1=\langle y\rangle$. The four generated-support groups are
$1,\langle y\rangle,\langle x\rangle,\langle x,y\rangle$ and satisfy all
intersection identities in (e). Their union generates only
$\langle x,y\rangle<G$. The coset geometry consists of two disjoint copies
of $K_{2,2}$, one in each coset of this subgroup, so it is not residually
connected. It is nevertheless thin and flag-transitive: every incident
pair has a common representative and is a translate of the base chamber.
Here the Borel is trivial. Thus triviality of the Borel alone does not
repair the claimed equivalence; the generation issue is independent.
\end{proof}
These examples concern the additional equivalences in Lemma 2.4, not the
separate residual-generation criterion in Lemma 2.3. They do not show that
every subsequent use fails; such a use must be checked with its actual
generation and faithfulness hypotheses.

\section{An outer automorphism quotient is not an action}
The opening of Section 5.2 (p.~18) takes $B=\Out(A)$ and says that it
naturally acts on $A$. The quotient $\Aut(A)/\Inn(A)$ does not in general
act on the elements of $A$ by choosing representatives: representatives
of one outer class differ by an inner automorphism. To obtain the intended
action inducing those outer classes, one needs a homomorphic section of
\[
 1\longrightarrow\Inn(A)\longrightarrow\Aut(A)
       \longrightarrow\Out(A)\longrightarrow1.
\]
This is more than a lack of a canonical choice.

\begin{proposition}\label{p:outer}
For the dihedral group
$A=\langle r,s\mid r^8=s^2=1,\ srs=r^{-1}\rangle$ of order $16$, the
map $\Aut(A)\to\Out(A)$ has no homomorphic section.
\end{proposition}
\begin{proof}
The elements of order eight generate the unique cyclic subgroup
$\langle r\rangle$, so every automorphism is uniquely
\[
 f_{a,b}: r\mapsto r^a,\quad s\mapsto r^b s,
 \qquad a\in\{1,3,5,7\},\quad b\in\mathbb Z/8.
\]
All these maps are automorphisms. Composition is
$f_{a,b}f_{c,d}=f_{ac,b+ad}$, with parameters modulo eight.
The inner automorphisms have $a\in\{1,7\}$ and even $b$; they are
generated by $f_{7,0}$ and $f_{1,2}$. Thus $\Out(A)\cong C_2^2$, with
generating classes $X=[f_{1,1}]$ and $Y=[f_{3,0}]$.

An involutory lift of $X$ has the form $f_{7,c}$ with $c$ odd: the other
possibility $f_{1,c}$ has square $f_{1,2c}\ne1$. An involutory lift of
$Y$ is either $f_{3,b}$ with $b$ even, or $f_{5,b}$ with $b\in\{0,4\}$.
Commuting $f_{7,c}$ with the first kind would require
$c+7b=b+3c$, or $c-3b=0\pmod4$, impossible by parity.
Commuting it with the second kind would require
$4c-6b=0\pmod8$, also impossible for odd $c$ and $b=0,4$.
A section would send the two commuting involutions $X,Y$ to commuting
involutory lifts, so no section exists.
\end{proof}
A chosen subgroup of actual automorphisms, such as the particular reversal
$\langle\tau\rangle$ in Theorem 5.2, does define an action; this objection
does not apply to that specified action. Nor does it deny the existence
of unrelated homomorphisms from $\Out(A)$ to $\Aut(A)$, such as the trivial
homomorphism. The missing interface is an action lifting the quotient
classes.

\section{Evidence and review boundary}
The accompanying standard-library program \texttt{finite\_models.py}
constructs left cosets and group operations independently of the author's
GAP script. It checks the $48$-element wreath model, every coordinate
support intersection and partial-flag extension for ranks three and four,
all triples of distinct involutions in the two reversal extensions, both
rank-two parabolic examples, and all involutory lifts used in
Proposition~\ref{p:outer}. The full arguments above establish the uniform
claims; the finite program also makes each particular countermodel
reproducible. No Lean formalization or whole-source recovery is claimed.
The three main proposed obstructions originated in the current N006
recovery work; this independent review verifies their source hypotheses
and supplies the direct involution proof and the additional preliminary
countermodels. No historical-priority claim is made.

\begin{thebibliography}{1}
\raggedright
\bibitem{source} C.~A.~Piedade and P.~Tranchida,
\emph{The Geometry of Wreath and Semi-Direct Products},
arXiv:2605.04646v2, submitted 28 May 2026; dated 29 May 2026, 22 pages.
\end{thebibliography}
\end{document}
